\subsection{B5. Holografie und Schatten: die konkrete native
Verbindung}\label{b5.-holografie-und-schatten-die-konkrete-native-verbindung}

\subsubsection{B5.1 Ein Kodierer ist bereits
vorhanden}\label{b5.1-ein-kodierer-ist-bereits-vorhanden}

Der vorhandene, frisch geprüfte Tensor erfüllt WW dagger gleich 8 I60.
Daher ist

\[\adjustbox{max width=\linewidth}{$\displaystyle 
V=W^\dagger/\sqrt8:\ \mathbb C^{60}\longrightarrow\Lambda^2\mathbb C^{64},
\qquad V^\dagger V=I_{60},\quad
\Pi=VV^\dagger=W^\dagger W/8
$}\]

eine exakte Isometrie mit einem Rang-60-Projektor Pi. Ein
60-dimensionaler logischer Zustand kann verlustfrei als kohärente
Überlagerung von Fermionpaaren dargestellt werden. Die andere Richtung
sieht aber nicht den ganzen 2016-dimensionalen Paarraum: Ihr Kern hat
Dimension 1956. Diese „dunklen`` Richtungen sind nicht deshalb generell
unphysikalisch; sie sind für genau diese Abbildung unsichtbar.

Die beiden Ansichten sind durch W also bereits verbunden. Die native
N=2-Dynamik im hellen Paar-plus-Boson-Bereich ist sogar exakt

\[\adjustbox{max width=\linewidth}{$\displaystyle 
H_{N=2,\mathrm{hell}}=
\begin{pmatrix}0&\sqrt8 g\\\sqrt8 g&\Delta\end{pmatrix}\otimes I_{60}.
$}\]

Sie wandelt die Kodierungsform um; am Prüfpunkt erreicht die
entsprechende Übergangswahrscheinlichkeit höchstens 2/27. Der nur in die
Paarseite eingebettete Code ist bei g ungleich null kein invarianter
H-Unterraum. Dieser kleine gemeinsame dynamische Träger ist ein
positiver Befund. Er ist N=2, nicht der native N=64-Grundzustand und
nicht dessen geladener Pol. Insbesondere ist er noch kein Transport
durch physikalischen Raum.

\subsubsection{B5.2 Warum das noch kein holografischer
Fehlerkorrekturcode
ist}\label{b5.2-warum-das-noch-kein-holografischer-fehlerkorrekturcode-ist}

Für jeden der 64 Fermionmoden wurde die logische Kompression seiner
Besetzungsmessung berechnet:

\[\adjustbox{max width=\linewidth}{$\displaystyle 
V^\dagger n_rV\quad\hbox{hat Spektrum}\quad
0\ (45\text{-fach}),\quad 1/8\ (15\text{-fach}).
$}\]

Das Resultat ist nicht skalar. Die Umgebung kann daher schon durch
Ablesen einer einzigen verlorenen Mode etwas über den logischen Zustand
erfahren. Die notwendige Fehlerkorrekturbedingung für die vollständige
60er-Codemenge ist verletzt: Der Code korrigiert nicht die beliebige
Löschung dieser einen Mode. Fermionparität ändert dieses
Besetzungsargument nicht, denn n\_r ist gerade. Kleinere geeignete
Untercodes oder ein aus nativen Operationen erzeugtes größeres
Kodierungsnetz sind dadurch nicht ausgeschlossen, aber noch nicht
konstruiert. Bloßes Aneinanderhängen desselben Tensors garantiert keinen
Code.

In etablierten holografischen Spielzeugmodellen sind isometrische
Kodierung, Rekonstruktion logischer Operationen aus verschiedenen
Randteilen und Fehlerkorrektur genau kontrollierte Eigenschaften. Das
ist ein hilfreiches Prüfschema, keine direkte Identifikation mit TFPT.
Siehe \href{https://arxiv.org/abs/1503.06237}{Pastawski, Yoshida, Harlow
und Preskill}. Die Rekonstruktion aus einem Teilrand ist auch in der
AdS/CFT-Arbeit von \href{https://arxiv.org/abs/1601.05416}{Dong, Harlow
und Wall} an konkrete Quanteninformations- und geometrische
Voraussetzungen gebunden. Diese Voraussetzungen wurden für TFPT nicht
nachgewiesen.

\subsubsection{B5.3 Was mehrere Schatten gemeinsam eindeutig machen
können}\label{b5.3-was-mehrere-schatten-gemeinsam-eindeutig-machen-kuxf6nnen}

Seien Ri festgelegte Mess- oder Reduktionsabbildungen eines gemeinsamen
Zustands rho. Für eine endliche, nicht eingeschränkte Zustandsklasse
sind alle Zustände aus den Schatten Ri(rho) genau dann unterscheidbar,
wenn auf den hermiteschen spurlosen Differenzen gilt

\[\adjustbox{max width=\linewidth}{$\displaystyle 
\bigcap_i\ker R_i=\{0\}.
$}\]

Das ist das einfache gemeinsame-Kern-Kriterium. Anschaulich: Was eine
Ansicht nicht sieht, muss eine andere sehen. Zwei Qubitansichten auf
(x,z) und (y,z) genügen gemeinsam zur Bloch-Rekonstruktion, jede für
sich nicht. Zweimal dieselbe Ansicht auf (x,z) genügt weiterhin nicht.
Das gilt ebenso für zwei Theorieberichte, deren Übereinstimmung bereits
aus denselben eingesetzten Annahmen stammt: Sie liefern keine zweite
unabhängige Messung.

Drei exakt nachgerechnete Grenzen sind wesentlich:

\begin{enumerate}
\def\labelenumi{\arabic{enumi}.}
\tightlist
\item
  Die orthogonalen Zustände (000 plus 111)/Wurzel 2 und (000 minus
  111)/Wurzel 2 haben dieselben reduzierten Zustände auf sämtlichen
  echten Teilmengen der drei Qubits. Eine gemeinsame phasensensitive
  XXX-Messung unterscheidet sie mit Erwartungswert +1 oder -1. Auch sehr
  viele lokale Schatten können also eine globale Phase übersehen.
\item
  Paarweise passende Überschneidungen garantieren nicht einmal einen
  gemeinsamen Ursprung. Drei binäre Variablen können nicht paarweise
  ausnahmslos verschieden sein. Jede einzelne perfekte Antikorrelation
  besitzt aber dieselben gleichverteilten Einzelmarginalen. Gemeint ist
  hier ein gemeinsames klassisches Tripel, nicht eine Behauptung gegen
  kontextuelle Quantenexperimente.
\item
  Selbst die vollständige Antwort einer Referenz bestimmt keinen völlig
  entkoppelten dunklen Zusatzsektor. H und H direkt plus K besitzen vom
  Startvektor (F,0) aus dieselben Energiemomente und dieselbe
  Zeitantwort. Ein exaktes kleines Matrixbeispiel prüft diesen
  Sachverhalt.
\end{enumerate}

Das erreichbare Ziel ist daher zunächst der \textbf{kleinste gemeinsame
beobachtbare Prozess}, nicht ein aus endlichen Schatten bewiesener
einzigartiger ontologischer Innenraum. Zwei Historien gelten operational
als gleich, wenn alle verfügbaren künftigen Eingriffs- und Auslesefolgen
dieselben Wahrscheinlichkeiten liefern. Das ist eine eindeutige
Unterscheidungsregel für gegebenes Verhalten; sie garantiert weder eine
eindeutige Hilbertraumdarstellung noch von selbst eine einfache
Geometrie.

\subsubsection{B5.4 Die drei Wände müssen unterschiedlich geprüft
werden}\label{b5.4-die-drei-wuxe4nde-muxfcssen-unterschiedlich-gepruxfcft-werden}

{\def\LTcaptype{none} % do not increment counter
\begin{longtable}[]{@{}
  >{\raggedright\arraybackslash}p{(\linewidth - 4\tabcolsep) * \real{0.3333}}
  >{\raggedright\arraybackslash}p{(\linewidth - 4\tabcolsep) * \real{0.3333}}
  >{\raggedright\arraybackslash}p{(\linewidth - 4\tabcolsep) * \real{0.3333}}@{}}
\toprule\noalign{}
\begin{minipage}[b]{\linewidth}\raggedright
Gemeinte Wand
\end{minipage} & \begin{minipage}[b]{\linewidth}\raggedright
Was sie gegenwärtig bedeutet
\end{minipage} & \begin{minipage}[b]{\linewidth}\raggedright
Welcher Anschluss fehlt
\end{minipage} \\
\midrule\noalign{}
\endhead
\bottomrule\noalign{}
\endlastfoot
Universalraum & Grenze zwischen gemeinsamem Operatorobjekt und seinen
reduzierten Ansichten & Ein gemeinsamer verfügbarer Operationssatz,
Zustand und Rekonstruktionsvertrag \\
TFPT & Grenze zwischen interner Compilerstruktur und physikalischer
Realisierung & Herleitung der Ausführung, relativer Parameter und
operationaler Teilung \\
Beobachtete Realität & Endlicher Zugang über Messungen, Präparationen
und Korrelationen & Ein quantitatives, nicht nachträglich angepasstes
Feld- und Messwörterbuch \\
\end{longtable}
}

Diese Erkenntnisgrenzen sind nicht schon drei bewiesene physikalische
Holografieschirme. Ein kosmologischer Horizont, ein Code-Rand und eine
offene mathematische Beweispflicht sind verschiedene Dinge. Ebenso sind
Universalraum und TFPT Kandidatenbeschreibungen der Realität, nicht
bereits drei experimentell etablierte wechselwirkende Welten.

Eine stärkere holografische Deutung müsste eine einzige Kodierung mit
Zuständen, Operationen und Zeitentwicklung verbinden. Auf einem
invarianten Codesektor wäre etwa V dagger V=I und Hrand V=V Hin eine
passende Intertwining-Bedingung. Bei nichtinvarianten reduzierten
Ansichten muss die entstehende Erinnerung mitgeführt werden; autonome
reduzierte Dynamik darf nicht einfach vorausgesetzt werden. Zusätzlich
braucht es eine Bedeutung von Randteilen, kontrollierte Rekonstruktion
und eine skalierende Geometrie.

\subsection{B6. Die einfachste neue Arbeitsrichtung und ihre
Abnahmekriterien}\label{b6.-die-einfachste-neue-arbeitsrichtung-und-ihre-abnahmekriterien}

Der plausible übersehene Punkt ist kein weiterer großer Tensor:
\textbf{Wir brauchen die gemeinsame Regel, die festlegt, welche
Operationen, Teilchenvariablen und Ansichten tatsächlich dieselbe
Ausführung beschreiben.} Noethers Symmetrieprüfung sagt, was erhalten
bleibt. Die Schattenprüfung sagt, was überhaupt unterscheidbar ist. Der
Eingriffstest sagt, was etwas anderes beeinflussen kann. Zusammen sind
diese drei kleinen Fragen strenger und informativer als weitere
isolierte Zahlenübereinstimmungen.

\begin{enumerate}
\def\labelenumi{\arabic{enumi}.}
\tightlist
\item
  \textbf{Den minimalen Ursprungsvertrag auswählen.} Den tatsächlichen
  Compiler auf den nativen kubischen Term, den zweiten kubischen Kanal
  und den quadratischen Paarterm prüfen. Die vollständige kleine Familie
  und der Bogoliubov-Test verhindern, dass derselbe Prozess mehrfach
  gezählt wird. Erfolg ist eine Ableitung oder ein präziser Ausschluss
  samt relativen Parametern, nicht nur die Feststellung erlaubter
  Symmetrie.
\item
  \textbf{Die Schattenkarte am vorhandenen W aufbauen.} Für denselben
  Vertrag Präparationen, Aufzeichnungen und kontrollierte
  Mehrzeitantworten in den verschiedenen Ansichten auflisten. Gemeinsam
  unsichtbare Richtungen und unverträgliche Überlappungen berechnen. Aus
  einem Teil der Daten eine Aussage bestimmen, die eine andere Ansicht
  ohne Nachjustierung testet. Der jetzt bewiesene Rang-60-Code und sein
  Einmoden-Leck sind Ausgangsdaten, kein Anlass, Raum oder perfekte
  Fehlerkorrektur bereits anzunehmen.
\item
  \textbf{Daraus einen gemeinsamen kausalen Teilträger gewinnen.} Eine
  verfügbare Operation muss die unbedingte Statistik eines anderen
  operational bestimmten Teils verändern. Derselbe Träger muss das
  relativistische Feldwörterbuch tragen. Erst dann sind Skalierung,
  chirales Maß und ein Spin-2-Sektor belastbare nächste Rechnungen.
\end{enumerate}

Die lokale Lanczos-Fortsetzung bleibt als kontrolliertes
Diagnoseinstrument nützlich, steht aber nicht mehr an erster Stelle der
fundamentalen Suche. Die drei Schritte ersetzen keine einzelnen
T1-T8-Nachweise. Auch eine erfolgreiche endliche Kodierung würde RH,
effiziente Faktorisierung oder P-versus-NP nicht automatisch lösen.

Der RH-Absatz der Eingabe enthält keine neue vollständige arithmetische
Beweiskonstruktion. Das vorhandene Forschungsregister wurde begrenzt
abgeglichen; sein Aktualitätslauf scheitert derzeit an einer nicht
verfügbaren historischen Quelle beziehungsweise Quellen-/Review-Drift.
Bestehende endliche Determinantsonden sind ausdrücklich Diagnosemodelle.
Es wurde kein neuer RH-Beweis, kein umfassender RH-Neulauf und keine
Erneuerung der entsprechenden Vertrauensmarker behauptet.

\textbf{Bilanz:} Ein gemeinsamer Kodierungsbaustein ist schon da;
außerdem sind die minimalen symmetrieverträglichen Dynamiken und zwei
wichtige Verwechslungen jetzt wesentlich schärfer bestimmt. Die
vollständige universelle physikalische Lösung bleibt offen. Der nächste
Schritt ist kleiner und konkreter geworden: dieselbe Ausführung durch
mehrere kalibrierte, dynamisch konsistente Ansichten rekonstruieren und
testen.

\subsection{B7. Abgleich des zuletzt eingesandten Berichts
v1.6.3}\label{b7.-abgleich-des-zuletzt-eingesandten-berichts-v1.6.3}

Der sechste Text dieser Runde, „Operationssatz, nativer Grundzustand und
Feldwörterbuch``, enthält überwiegend bereits berücksichtigte Befunde:
die Kommutantenleiter, das Lochbild, die Normen der ersten Stufen, den
nichtverschwindenden Seitenzweig und den verschwindenden
Weyl-Skalarkanal. Seine zweite vorgeschlagene Folgeaufgabe wurde
inzwischen wesentlich weitergeführt: Alle vier Zwei-Boson-Singuletts
sind konstruiert und der richtige H-Lanczos-Anfang ist fünf Glieder weit
bestimmt.

Vier Präzisierungen verhindern eine Übernahme veralteter Schlüsse:

\begin{enumerate}
\def\labelenumi{\arabic{enumi}.}
\tightlist
\item
  \textbf{Sieben und sechzehn sind verschiedene Kommutanten.} Sieben
  gilt für volle innere Gruppe zusammen mit X und Nb; sechzehn für die
  Gruppe allein. Die drei hellen inneren Typen kommen auf beiden Seiten
  der Umwandlung vor. Einzelne Tabellenzeilen sind multiplizitätsfrei,
  der ganze N=3-Sektor aber nicht. Der Schluss vom nichtskalaren Clock
  auf eine äußere, nicht zusammenhängende Symmetrie ist weiterhin
  falsch.
\item
  \textbf{Eine kleine Zweignorm beweist keine globale Konvergenz.} Der
  genannte Defekt 2,90 mal 10 hoch -5 betrifft eine bestimmte
  Kontraktion. Die höhere Bosonstufe ist weiterhin gekoppelt. „Numerisch
  fast allein durch die Normen lösbar`` und „vernachlässigbar`` sind
  ohne Restschranke zu stark. Auch Zh plus Zadd gleich eins ist eine
  Summenregel, keine Bestimmung sämtlicher Spektrallinien.
  Ritz-Erwartungswerte bleiben Erwartungswerte des Ritz-Zustands, nicht
  automatisch von Omega.
\item
  \textbf{Der Feldschluss ist ansatzabhängig.} Der angegebene Fünfzyklus
  (16,45,0,30,37) wurde frisch am W überprüft. Damit ist eine einzige
  binäre Chiralitätsbelegung der primitiven Labels mit entgegengesetzter
  Händigkeit auf jeder Kante unmöglich. Das verbietet kein allgemeines
  Dirac-Wörterbuch mit erweiterten Trägern, Ableitungen oder
  zusammengesetzten Feldern. Der Tensorfeldkanal ist unter dem
  angegebenen engen Ansatz sinnvoll, nicht schon der universell einzige
  physikalische Feldtyp.
\item
  \textbf{Der gespeicherte Rechnerstand ist älter als der Text.} Die
  lokal zugehörige Datei \codeword{native\_ground\_state.json} hat einen
  anderen Checkerhash als das aktuelle Programm. Ihr gespeichertes
  w2-Normquadrat ist um 229 Quadrat gegenüber der physikalischen Norm
  skaliert. Das aktuelle Programm enthält bereits die richtige
  Entskalierung; die alte Datei ist also kein aktueller Replaynachweis.
  Sie meldet außerdem \codeword{norms\_by\_traces\_used=false} und
  Ritz-Tiefe drei. Die im Text erwähnte Datei
  \codeword{norms\_by\_traces.json} und das dortige Replaymanifest waren
  bei dieser Prüfung nicht vorhanden. Diese Beobachtung schließt ein
  Resultat an anderem Ort nicht aus, liefert hier aber keines.
\end{enumerate}

Daher wird insbesondere die Behauptung, nu4 sei in dieser Quelle auf
zwei unabhängigen Wegen vollständig reproduziert, nicht übernommen. Der
Text selbst bezeichnet nu4 zugleich als „nur Spurnetzwerk`` und
beschreibt den Brute-Force-Abgleich nur bis nu3. Der frühere gepinnte
exakte nu4-Eingang unserer eigenen Rechnung bleibt davon getrennt und
wurde nicht als neue Enumeration ausgegeben. Ein fertig berechnetes nu5
wird aus dem vorhandenen Programm mit \codeword{MAX\_N=5} nicht
abgeleitet; eine Einstellung ist kein Ergebnis.

\textbf{Die interessante verbleibende Richtung ist die
Spurnetzwerk-Methode.} Sie kann Normen durch Kontraktionen des
vorhandenen Tensors berechnen, statt alle Vielteilchenamplituden zu
speichern. Zur fundamentalen Vereinfachung wird sie erst dann, wenn sie
auch gemischte H-Momente, Seitenzweige und kontrollierte
Mehrzeitantworten liefert. Der konkrete nächste Test ist daher kein
bloßer weiterer Normwert: Eine zweite kleine Kontraktionsrechnung soll
die schon exakt bekannten zehn H-Momente reproduzieren und dann
mindestens ein noch nicht bekanntes gemischtes Matrixelement samt
Fehlerkontrolle bestimmen. Diese Methode lässt sich mit der
Schattenkarte verbinden, ersetzt aber keine native Operation.

Das Prüfpaket bewahrt neben dem Text die beiden widersprechenden
Programmstempel und das vorhandene Spurnetzwerkprogramm. Diese fremden
Dateien wurden nicht verändert oder als vollständig neu ausgeführt
ausgegeben.

\begin{center}\rule{0.5\linewidth}{0.5pt}\end{center}

\section{Historischer Anhang: vollständige Fassung
v1.6.6}\label{historischer-anhang-vollstuxe4ndige-fassung-v1.6.6}

Die folgende vollständige Fassung wird unverändert bewahrt. Maßgeblich
für den aktuellen Status sind die vorangestellten Kapitel v1.6.7. Ältere
Vorhaben, Modellannahmen und Zahlen erhalten dadurch keinen neuen
Beweisstatus; die kleinen Energieverschärfungen und die neuen Singulett-
und Bilinearergebnisse stehen oben.

\section{TFPT / Universalraum: Clock, gemeinsame Quelle und
tatsächlicher
Transport}\label{tfpt-universalraum-clock-gemeinsame-quelle-und-tatsuxe4chlicher-transport}

\subsection{Konsolidierung und eigene Forschungsfortsetzung v1.6.6 - 15.
September
2026}\label{konsolidierung-und-eigene-forschungsfortsetzung-v1.6.6---15.-september-2026}

\textbf{Ergebnis:} Der dokumentierte endliche Clock ist jetzt
ausdrücklich als Element der vorhandenen inneren Spin(10)-Symmetrie
konstruiert. Dadurch lässt sich die bislang offene gemeinsame
Clock-/Casimir-Auslesung exakt bestimmen. Daneben wurde ein kleiner,
echter Vierzustandskanal der ursprünglichen Paarwechselwirkung gefunden,
der nach Zulassung eines Bosonmischers eine Fermionenmarke überträgt.
Die Herkunft dieses Mischers ist nicht bewiesen.

Der neue Universalraum-Text enthält eine sinnvolle Suchrichtung -
gemeinsame Quelle statt vorab unabhängiger Banken -, aber Überlappung
allein erzeugt weder Transport noch Raumzeit, Eichkrümmung oder die
arithmetische Spur. Mehrere seiner vermeintlich automatischen Übergänge
werden hier präzisiert.

Diese Revision integriert \textbf{drei neue Nutzertexte}, die
unabhängigen parallelen Prüfungen und eigene neue Herleitungen. Die
vollständige frühere Konsolidierung v1.6.5 einschließlich v1.6.4 bleibt
im historischen Anhang erhalten. Das ist eine aktualisierte
Forschungsfassung in Markdown mit kurzem Änderungsdokument und einfacher
Erklärung; keine behauptete neue PDF-/Web-Publikation oder vollständige
Theory of Everything.

